\textbf{结论名称：} 隐零点问题三步法（设零点 $\to$ 代换 $\to$ 消去零点）

\textbf{适用条件：} 导数 $f'(x)$ 的零点不可显式求解，但可判断存在且唯一；函数含 $e^x$ 与 $\ln(x+a)$（$a$ 为常数）等超越结构，零点方程能提供指数与对数的互化关系。

\textbf{变量说明：} 设 $x_0$ 为 $f'(x)$ 的隐零点，常数 $a$ 满足 $x_0+a>0$，由 $f'(x_0)=0$ 整理得 $e^{x_0}(x_0+a)=1$。

\textbf{核心公式：}
\[
\boxed{ e^{x_0}(x_0+a)=1 \quad\Longleftrightarrow\quad \ln(x_0+a)=-x_0 }
\]

\textbf{使用场景：} 求含指数与对数的函数最值或证明不等式时，利用隐零点关系将超越项整体替换，转化为多项式或有理分式，进而用基本不等式或二次函数处理。

\textbf{简单例子：} 证明 $f(x)=e^x-\ln x>2$。$f'(x)=e^x-\frac{1}{x}$ 在 $(0,+\infty)$ 上存在唯一零点 $x_0\in(0,1)$（$a=0$），满足 $e^{x_0}=\frac{1}{x_0}$。代换得 $\ln x_0=-x_0$，于是
\[
f(x_0)=e^{x_0}-\ln x_0=\frac{1}{x_0}+x_0>2\quad (x_0\neq 1),
\]
其中 $f(x_0)$ 为最小值，故恒有 $f(x)>2$。